\section{Intervention and Participant Discontinuation/Withdrawal}

This section defines the prospectively specified rules for discontinuing the study
intervention (the device, the drug, or both), the procedures governing a
participant's withdrawal from the study, and the handling of participants lost to
follow-up. Discontinuation of an intervention is distinct from withdrawal from the
study: a participant who stops one or both components for a protocol-defined reason
remains in the study for safety and survival follow-up unless consent for that
follow-up is withdrawn.

\subsection{Discontinuation of Study Intervention}

\paragraph{Device stopping rules.}
The investigational device is discontinued for the index procedure, with conversion
to the appropriate fallback technique, whenever any of the following prespecified
conditions occurs. These rules implement the Physical AI stopping criteria of 21 CFR
$\S$312.23(g)(i)(v) and the divergence-reporting threshold of $\S$312.32(g)(v)
\cite{kawchak2026adapt312}.

\begin{itemize}[leftmargin=1.4em,itemsep=1pt,topsep=2pt]
\item \textit{Conversion to open or manual technique.} Any intra-operative judgment
by the supervising surgeon that patient safety or oncologic adequacy is better
served by manual robotic or open conversion (including uncontrolled bleeding,
unfavorable anatomy, or loss of a critical view) ends device-directed operation
immediately; conversion is always available and is never itself a protocol
violation.
\item \textit{Repeated force-limit excursions.} Repeated excursions against the
$\leq 3$ N per-arm or $\leq 18$ N cumulative cross-arm tip-force caps, or any
sustained or consecutive pattern of force-clamp activations indicating that the
case cannot proceed within the force envelope, trigger discontinuation.
\item \textit{No-fly-zone violations.} Any breach of a vascular no-fly zone around a
named vessel (superior mesenteric vein or portal vein at 1.0 mm; hepatic artery,
celiac axis, or superior mesenteric artery at 1.5 mm), or a recurring pattern of
no-fly approaches, ends device-directed dissection in that field.
\item \textit{Sim-to-real divergence beyond twice tolerance.} Any instance in which
the observed device behavior diverges from the validated simulation prediction by
more than twice the sim-to-real gap tolerance (a trajectory deviation greater than
4 mm or a force discrepancy greater than 1.0 N) triggers discontinuation and a
report under $\S$8.3.6.
\item \textit{Hard-stop E-stop event.} Any hard-stop cross-arm E-stop (the
$\leq 3$ ms full-platform halt, whether software- or hardware-initiated) ends the
device-directed segment; resumption with the device requires that the cause be
identified and that the system re-pass the pre-procedure safety matrix
($\S$6.2), failing which the case is completed by conversion.
\end{itemize}

Following any device discontinuation, the operation is completed by the most
appropriate conventional means, the hash-chained audit trail is preserved across
the event window ($\S$8.3.6), and the event is reviewed by the Physical AI Safety
Review Committee and the data safety monitoring board before the device is used in
a subsequent case.

\paragraph{Drug stopping rules.}
Daraxonrasib is discontinued in an individual participant for a dose-limiting
toxicity attributable to the drug during the 28-day DLT window, for any other
clinically significant or unacceptable toxicity that does not resolve with protocol
dose interruption and reduction, for a confirmed ISGPS Grade C postoperative
pancreatic fistula \cite{Bassi2017} (which precludes restart and maps to the T+21d
or continued-hold arm of the perioperative advisory, $\S$6.1.2), for pregnancy, for
intercurrent illness that contraindicates continuation, or at the participant's
request. Dose interruptions and reductions for lower-grade toxicity follow the
protocol algorithm and do not by themselves constitute permanent discontinuation;
a participant who permanently discontinues the drug remains in the study for safety
and survival follow-up.

\subsection{Participant Discontinuation/Withdrawal from the Study}

A participant may withdraw from the study at any time, for any reason, without
penalty and without loss of any benefit to which the participant is otherwise
entitled, and a participant may separately withdraw consent for continued
study-specific procedures while permitting ongoing survival follow-up through the
medical record. The investigator may withdraw a participant for safety, for
intercurrent illness, for significant non-compliance that compromises participant
safety or data integrity, for loss of eligibility, or at the sponsor's direction if
the study is terminated. The reason for and date of any discontinuation or
withdrawal are recorded in the case report form, and a participant who withdraws is
asked to complete a final safety assessment where feasible.

Data and specimens collected up to the point of withdrawal are retained and analyzed
in accordance with the consent and applicable regulation; withdrawal of consent for
future participation does not require deletion of data already collected. The
relationship of discontinuation and withdrawal to the analysis populations (the
safety, DLT-evaluable, and modified intention-to-treat populations) is defined in
the statistical considerations ($\S$9). A participant who withdraws before
completing the operative intervention is replaced for the purpose of completing the
3+3 cohort only if not DLT-evaluable, per the rules in $\S$9.

\subsection{Lost to Follow-Up}

A participant is considered lost to follow-up only after the site has made and
documented at least three contact attempts by telephone across separate days,
followed by a written contact attempt to the last known address, over the 24-month
overall-survival window. Each attempt is recorded with its date and outcome in the
case report form. Before declaring a participant lost to follow-up, the site
attempts, with appropriate authorization, to ascertain vital status from the medical
record, the referring provider, and permissible public sources, because survival is
a key endpoint of the study. A participant who cannot be reached after these
attempts is classified as lost to follow-up, with the last date known alive
recorded. The handling of the resulting missing data, including censoring rules for
progression-free and overall survival and the sensitivity analyses used to assess
the impact of missingness, is specified in $\S$9.4.
